\chapter{Visão Lógica (Logical View)}
\label{cap_visao_logica}

A Visão Lógica do Modelo 4+1 descreve a organização conceitual e estrutural do software, identificando as abstrações do domínio, os subsistemas lógicos, as camadas de software e os padrões de projeto que regem a decomposição funcional do sistema \cite{kruchten1995,martin2018}. No SIGEA-GTT, a Visão Lógica fundamenta-se nos princípios de \textit{Clean Architecture} e \textit{Domain-Driven Design} (DDD) \cite{evans2003}, assegurando baixo acoplamento entre a lógica clínica de negócio e os mecanismos de infraestrutura técnica.

% ===========================================================================
\section{Decomposição em Camadas do Backend (Spring Boot 3.3)}
\label{sec_camadas_backend}
% ===========================================================================

O backend foi implementado em Java 21 LTS sob o framework Spring Boot 3.3.x, organizado em quatro camadas lógicas horizontais estritamente delimitadas (\autoref{fig_camadas_backend}):

\begin{figure}[htb]
\centering
\caption{Arquitetura em Camadas do Backend SIGEA-GTT.}
\label{fig_camadas_backend}
\resizebox{0.85\textwidth}{!}{%
\begin{tikzpicture}[
  >=Stealth,
  layer/.style={
    rectangle, draw=blue!80!black, fill=blue!8, thick,
    minimum width=11cm, minimum height=1.3cm,
    font=\small\bfseries, align=center, rounded corners=4pt
  },
  note/.style={font=\scriptsize, text=gray!80!black}
]

\node[layer, fill=blue!15] (pres) at (0, 4.5) {
  Camada de Apresentação (REST Controllers \& OpenAPI 3)\\
  \textnormal{\scriptsize Controladores REST, Bean Validation, DTOs, Mappers MapStruct, RestControllerAdvice}
};

\node[layer, fill=blue!10] (app)  at (0, 2.7) {
  Camada de Aplicação e Serviços (@Service)\\
  \textnormal{\scriptsize Orquestração de Negócio, Transações @Transactional, Regras IHI-GTT e Cálculo Epidemiológico}
};

\node[layer, fill=blue!6]  (dom)  at (0, 0.9) {
  Camada de Domínio (Domain Entities \& Value Objects)\\
  \textnormal{\scriptsize Entidades JPA/Hibernate, Enums de Negócio, Modelos de Prontuário e Ferramentas JSONB}
};

\node[layer, fill=blue!3]  (infra)at (0, -0.9) {
  Camada de Infraestrutura e Persistência\\
  \textnormal{\scriptsize Repositórios Spring Data JPA, Migrations Flyway, Spring Security 6 (JWT) e Driver PostgreSQL}
};

\draw[->, very thick, blue!80!black] (pres.south) -- (app.north);
\draw[->, very thick, blue!80!black] (app.south)  -- (dom.north);
\draw[->, very thick, blue!80!black] (dom.south)  -- (infra.north);

\end{tikzpicture}%
}
\fonte{Elaborado pelo autor com TikZ (2026).}
\end{figure}

\begin{enumerate}
    \item \textbf{Camada de Apresentação (Controllers REST)}:
    \begin{itemize}
        \item Expõe endpoints RESTful padronizados sob rotas semânticas (ex: \texttt{/api/v1/turmas}, \texttt{/api/v1/atividades}, \texttt{/api/v1/submissoes});
        \item Recebe e retorna exclusivamente Objetos de Transferência de Dados (\textit{Data Transfer Objects} -- DTOs), impedindo a exposição direta de entidades JPA na fronteira HTTP;
        \item Aplica validação declarativa de entrada via \textit{Bean Validation} (\texttt{@Valid}, \texttt{@NotBlank}, \texttt{@Pattern});
        \item O tratamento de exceções é centralizado com \texttt{@RestControllerAdvice}, mapeando falhas de negócio em respostas JSON padronizadas (\textit{Problem Details} -- RFC 7807);
        \item Autodocumentação síncrona gerada pelo SpringDoc OpenAPI 3 (\path{/swagger-ui.html}).
    \end{itemize}

    \item \textbf{Camada de Aplicação / Serviço (\texttt{@Service})}:
    \begin{itemize}
        \item Encapsula as regras de negócio clínicas e acadêmicas, controlando o fluxo transacional com a anotação \texttt{@Transactional(rollbackFor = Exception.class)};
        \item Executa os cálculos dos indicadores epidemiológicos oficiais do IHI ($TI_{1000}$, $TI_{100}$, $P_{EA}$);
        \item Converte entidades em DTOs e vice-versa de forma imutável utilizando mappers compilados pelo MapStruct (evitando custos computacionais de reflexão em tempo de execução).
    \end{itemize}

    \item \textbf{Camada de Domínio}:
    \begin{itemize}
        \item Contém as entidades relacionais gerenciadas pelo Hibernate 6, anotadas com Lombok (\texttt{@Getter}, \texttt{@Setter}, \texttt{@NoArgsConstructor}) para eliminação de código redundante;
        \item Implementa os enums que governam os estados do sistema: \texttt{UserRole} (\texttt{ADMIN}, \texttt{PROFESSOR}, \texttt{STUDENT}), \texttt{TurmaStatus} (\texttt{ABERTA}, \texttt{ENCERRADA}) e \texttt{SubmissaoStatus} (\texttt{SUBMETIDO}, \texttt{CORRIGIDO});
        \item Suporte nativo a tipos complexos \texttt{JSONB} do PostgreSQL 16 para Prontuários Clínicos e Ferramentas da Qualidade.
    \end{itemize}

    \item \textbf{Camada de Infraestrutura e Persistência}:
    \begin{itemize}
        \item Provedores Spring Data JPA estendendo \texttt{JpaRepository<T, ID>}, com consultas otimizadas através de \textit{derived query methods} e JPQL indexadas;
        \item Versionamento evolucionário do esquema do banco de dados relacional via migrações Flyway (\texttt{db/migration/V1\_\_...});
        \item Camada de segurança configurada pelo Spring Security 6 com filtro customizado \texttt{JwtAuthenticationFilter}, gerenciando tokens \textit{bearer stateless}.
    \end{itemize}
\end{enumerate}

% ===========================================================================
\section{Decomposição Modular do Frontend (Angular 18+)}
\label{sec_camadas_frontend}
% ===========================================================================

O frontend adota a arquitetura moderna do Angular 18+, estruturado inteiramente com \textbf{Standalone Components} (eliminando o conceito legado de \texttt{NgModule}) e reatividade baseada em \textbf{Angular Signals}. A base de código organiza-se em três macro-camadas:

\begin{enumerate}
    \item \textbf{Camada de Núcleo (\texttt{core/})}:
    \begin{itemize}
        \item Serviços globais \textit{singleton}: autenticação (\texttt{AuthService}), persistência de tokens (\texttt{Token\-Storage\-Service}) e notificações (\texttt{Notification\-Service});
        \item Interceptores HTTP globais: \texttt{JwtInterceptor} (injeta o token JWT no cabeçalho das requisições) e \texttt{LoadingInterceptor} (gerencia o estado reativo de carregamento global com ativação de \textit{spinner} sobreposto em chamadas assíncronas);
        \item Guardas de rota funcionais (\textit{Functional Guards}): \texttt{authGuard} (restringe rotas privadas a usuários autenticados), \texttt{roleGuard} (valida perfis de acesso) e \texttt{firstAccessGuard} (bloqueia o uso do sistema caso a troca obrigatória de senha esteja pendente).
    \end{itemize}

    \item \textbf{Camada de Funcionalidades (\texttt{features/})}:
    \begin{itemize}
        \item Módulos de tela funcionais: \texttt{auth}, \texttt{admin}, \texttt{turmas}, \texttt{atividades}, \texttt{auditoria-gtt}, \texttt{ferramentas}, \texttt{correcao} e \texttt{dashboard};
        \item Cada componente de funcionalidade consome serviços de aplicação dedicados e formulários reativos fortemente tipados (\texttt{FormGroup}, \texttt{FormControl}).
    \end{itemize}

    \item \textbf{Camada Compartilhada (\texttt{shared/})}:
    \begin{itemize}
        \item Componentes visuais reutilizáveis: modais genéricos de confirmação de duas etapas para prevenção de ações destrutivas (\texttt{RNF08}), cabeçalhos de tabela com ordenação, controles de paginação fixos de exatamente 10 itens (\texttt{RNF05}) e barras de progresso cromáticas hospitalares;
        \item \textit{Pipes} personalizados de formatação: formatação de período acadêmico, renderização de scores GUT e máscaras de data/hora.
    \end{itemize}
\end{enumerate}

% ===========================================================================
\section{Diagrama de Classes de Domínio}
\label{sec_diagrama_classes_dominio}
% ===========================================================================

O Diagrama de Classes de Domínio da \autoref{fig_classes_dominio} expressa as entidades JPA fundamentais do SIGEA-GTT, seus atributos centrais e as associações estruturais que espelham as regras de negócio clínicas e educacionais.

\begin{figure}[htb]
\centering
\caption{Diagrama de Classes de Domínio do SIGEA-GTT (UML 2.5).}
\label{fig_classes_dominio}
\resizebox{0.95\textwidth}{!}{%
\tikzset{
  cls/.style={
    rectangle, draw=blue!70!black, fill=white, thick,
    minimum width=3.8cm, font=\scriptsize, align=left
  },
  clshead/.style={
    rectangle, draw=blue!70!black, fill=blue!15, thick,
    minimum width=3.8cm, minimum height=0.65cm,
    font=\small\bfseries, align=center
  }
}
\begin{tikzpicture}[>=Stealth, node distance=0.35cm]

% User
\node[clshead] (uh) at (0,0)   {User};
\node[cls, anchor=north] (u) at (uh.south)
  {id: UUID\\email: String\\senha: String\\nome: String\\role: UserRole\\matricula: String\\primeiroAcesso: Boolean};

% Turma
\node[clshead] (th) at (5.5,0)   {Turma};
\node[cls, anchor=north] (t) at (th.south)
  {id: Long\\nome: String\\materia: String\\periodo: String\\status: TurmaStatus\\professor: User};

% Atividade
\node[clshead] (ah) at (11.5,0)  {Atividade};
\node[cls, anchor=north] (a) at (ah.south)
  {id: Long\\titulo: String\\descricao: String\\prazoEntrega: LocalDateTime\\prontuarioClinico: JSONB};

% Submissao
\node[clshead] (sh) at (11.5,-4.5) {Submissao};
\node[cls, anchor=north] (s) at (sh.south)
  {id: Long\\dataSubmissao: LocalDateTime\\tempoGastoSegundos: Integer\\nota: BigDecimal\\parecerFormativo: String\\status: SubmissaoStatus\\ferramentasQualidade: JSONB};

% ModuloGTT
\node[clshead] (mh) at (0,-4.5)  {ModuloGTT};
\node[cls, anchor=north] (m) at (mh.south)
  {id: Long\\codigo: String\\nome: String\\descricao: String};

% GatilhoGTT
\node[clshead] (gh) at (5.5,-4.5)  {GatilhoGTT};
\node[cls, anchor=north] (g) at (gh.south)
  {id: Long\\codigo: String\\descricao: String\\pistaClinica: String};

% GravidadeMERP
\node[clshead] (rh) at (0,-8.5)  {GravidadeMERP};
\node[cls, anchor=north] (r) at (rh.south)
  {id: Long\\categoria: String\\danoReal: Boolean\\descricao: String};

% SubmissaoGatilho
\node[clshead] (sgh) at (5.5,-8.5) {SubmissaoGatilho};
\node[cls, anchor=north] (sg) at (sgh.south)
  {id: Long\\danoReal: Boolean\\justificativaClinica: String};

% Relacionamentos
\draw[->, thick] (th.west) -- node[above, font=\tiny]{1 professor} (uh.east);
\draw[->, thick] (ah.west) -- node[above, font=\tiny]{N} (th.east);
\draw[->, thick] (sh.north) -- node[right, font=\tiny]{N} (ah.south);
\draw[->, thick] (gh.west) -- node[above, font=\tiny]{N} (mh.east);
\draw[->, thick] (sgh.north) -- node[right, font=\tiny]{N} (gh.south);
\draw[->, thick] (sgh.west) -- node[above, font=\tiny]{N} (r.east);
\draw[->, thick] (sgh.east) -- node[above, font=\tiny]{N:1} (s.south west);
\draw[->, thick] (s.west) -- node[above, font=\tiny]{aluno} (t.south east);

\end{tikzpicture}%
}
\fonte{Elaborado pelo autor com TikZ (2026).}
\end{figure}

% ===========================================================================
\section{Diagrama de Entidade-Relacionamento (DER)}
\label{sec_der}
% ===========================================================================

A persistência do SIGEA-GTT é gerenciada pelo banco de dados relacional PostgreSQL 16 LTS. A \autoref{fig_der_arquitetura} apresenta o Diagrama de Entidade-Relacionamento do sistema, evidenciando as chaves primárias, chaves estrangeiras e a modelagem híbrida relacional/JSONB.

\begin{figure}[htb]
\centering
\caption{Diagrama de Entidade-Relacionamento (DER) do SIGEA-GTT.}
\label{fig_der_arquitetura}
\resizebox{0.92\textwidth}{!}{%
\begin{tikzpicture}[
  >=Stealth,
  entity/.style={
    rectangle, draw=blue!70!black, fill=blue!8, thick,
    minimum width=3.2cm, minimum height=0.8cm,
    font=\small\bfseries, rounded corners=2pt
  },
  rela/.style={
    diamond, draw=orange!80!black, fill=yellow!10, thick,
    minimum width=2.4cm, minimum height=0.75cm,
    font=\tiny\bfseries, aspect=2.4
  },
  lbl/.style={font=\tiny, midway}
]

% Linha 1
\node[entity] (users)     at ( 0,  0) {users};
\node[rela]   (r1)        at ( 3.5, 0) {leciona};
\node[entity] (turmas)    at ( 7,  0) {turmas};
\node[rela]   (r2)        at (10.5, 0) {possui};
\node[entity] (atividades)at (14,  0) {atividades};

% Linha 2
\node[rela]   (r3)        at ( 3.5, -3) {matricula};
\node[rela]   (r4)        at (14, -2)   {gera};
\node[entity] (submissoes)at (14, -4.5) {submissoes};

% Linha 3
\node[entity] (modulos)   at ( 0, -5) {modulos\_gtt};
\node[rela]   (r5)        at ( 3.5, -5) {contem};
\node[entity] (gatilhos)  at ( 7, -5) {gatilhos\_gtt};
\node[rela]   (r6)        at (10.5, -5) {associa};

% Linha 4
\node[entity] (merp)      at ( 7, -8) {gravidades\_merp};
\node[rela]   (r7)        at (10.5, -8) {classifica};
\node[entity] (subgat)    at (14, -8) {submissoes\_gatilhos};

% Conexoes
\draw[-] (users) -- node[lbl, above]{1} (r1);
\draw[->](r1)    -- node[lbl, above]{N} (turmas);

\draw[-] (turmas)    -- node[lbl, above]{1} (r2);
\draw[->](r2)        -- node[lbl, above]{N} (atividades);

\draw[-] (turmas)    -- node[lbl, left]{N}  (r3);
\draw[->](r3)        -- node[lbl, left]{M} (users);

\draw[-] (atividades)-- node[lbl, right]{1} (r4);
\draw[->](r4)        -- node[lbl, right]{N} (submissoes);

\draw[-] (modulos)   -- node[lbl, above]{1} (r5);
\draw[->](r5)        -- node[lbl, above]{N} (gatilhos);

\draw[-] (gatilhos)  -- node[lbl, above]{1} (r6);
\draw[->](r6)        -- node[lbl, above]{N} (subgat);

\draw[-] (submissoes)-- node[lbl, right]{1} (subgat);
\draw[-] (merp)      -- node[lbl, above]{1} (r7);
\draw[->](r7)        -- node[lbl, above]{N} (subgat);

\end{tikzpicture}%
}
\fonte{Elaborado pelo autor com TikZ (2026).}
\end{figure}
